\section{Proof of Asymptotic Efficiency} \label{app_f: Asymp Efficiency}
This appendix includes the proof of Theorem~\ref{thm:Efficient Influence Curve Approximation} and Theorem~\ref{thm:plugin_HAL_asymptotic_efficiency}. We provide a standard TMLE proof establishing asymptotic efficiency under the assumption that $D^{(k)}_M({\cal R}_n)$ approximates $D^{(k)}_U([0,1])$. 

\subsection{Proof of Theorem~\ref{thm:Efficient Influence Curve Approximation}}
\begin{proof}
We have that $P_n S_{f_n}(\phi_j)=0$ for the $J_n - 1$ scores that define the $L_1$ constrained score space in $D^{(k)}({\cal R}_n)$. The tangent space generated by paths through $f$ is given by closure of linear span $\{S_f(\phi): \phi\in D^{(k)}_U([0,1])\}$. Thus, we can represent 
$D^*_f=S_f(\phi_f)$ for some $\phi_f\in D^{(k)}_U([0,1])$. Therefore, if the basis ${\cal R}_n$ is chosen so that $D^{(k)}({\cal R}_n)$  yields a $O^+(1/J_n^{k+1})$ $L^2$-approximation of $D^{(k)}_U([0,1])$ and this $L^2$ approximation can be extended to the corresponding $L_1$ constrained score space, then we can find a projection $\tilde{\phi}_f$ of $\phi_f$ onto the $L_1$ constrained score space of $D^{(k)}({\cal R}_n)$ that approximates $\phi_f$ in $L^2$-norm within distance $O^+(1/J_n^{k+1})$.
Therefore, following the same argument as in Appendix~\ref{App C: Score Analysis}, we have

\begin{align*}
P_n D^*_{f_n}=P_n S_{f_n}(\phi_{f_n})
&=P_n\{S_{f_n}(\phi_{f_n})-S_{f_n}(\tilde{\phi}_{f_n})\} + P_nS_{f_n}(\tilde{\phi}_{f_n})\\
&=(P_n-P_0)\{\phi_{f_n} - \tilde{\phi}_{f_n}\}+ \\
&\quad(P_0- P_{f_n})\{\phi_{f_n} - \tilde{\phi}_{f_n}\} + 0,
\end{align*}

where we note that the expectation of scores at $f_n$ under $P_{f_n}$ equals zero, because $S_{f_n}(\tilde{\phi}_{f_n}) = \tilde{\phi}_{f_n} - P_{f_n}\tilde{\phi}_{f_n}$ by definition.

\paragraph{Bounding the First Empirical Process Term}
We could directly use that the covering number of the basis class is the covering number of $D^{(k)}_U([0,1])$. And using the $L^2$ covering number from Appendix~\ref{appB: $L^2$ Convergence Analysis}, we can bound the first term can be bounded by $n^{-1/2}J_{2}(\delta_n,D^{(k)}_U([0,1]))$ with $\delta_n \asymp n^{-\frac{k+1}{2k+3}}$, where $J_{2}(\delta,D^{(k)}_U([0,1])\asymp\delta^{(2k+1)/(2k+2)}$. Thus, we have 
\begin{align*}
\left| (P_n - P_0)\left\{ \phi_{f_n} - \tilde{\phi}_{f_n} \right\} \right| 
&= O_P\left( n^{-1/2} J_2(\delta_n, D^{(k)}_U([0,1])) \right) \notag \\
&= O_P\left( n^{-1/2} \delta_n^{\frac{2k+1}{2k+2}} \right) \notag \\
&= O_P\left( n^{-\frac{2k+2}{2k+3}} \right) \notag \\
&= o_P\left( n^{-1/2} \right). 
\end{align*}

\paragraph{Bounding the Second Term}
The second term can be bounded with Cauchy-Schwarz inequality by $\| p_{f_n}-p_0\|_{\mu}\| \phi_{f_n}-\tilde{\phi}_{f_n}\|_{\mu}$, where we shorthanded $\|\cdot\|_{L^2(\mu)}$ by $\|\cdot\|_{\mu}$.
By the $L^2$ convergence Theorem~\ref{thm:l2-convergence-hal-mle}, we have $\| p_{f_n}-p_0\|_{\mu}=O_p(n^{-(k+1)/(2k+3)})$. By the $L^2$ Approximation Assumption, $\| \phi_{f_n}-\tilde{\phi}_{f_n}\|_{\mu}= O^+(J_n^{-(k+1)})$. So that we obtain the bound $O_p^+(n^{-(k+1)/(2k+3)} J_n^{-(k+1)})$. Choose $J_n \asymp^+ n^{1/(2k+3)}$, which could be achieved by cross validation or slight undersmoothening. Again, it follows that for all $k$, the second term is $o_P(n^{-1/2})$. 

\begin{align*}
\left| (P_0 - P_{f_n})\left\{ \phi_{f_n} - \tilde{\phi}_{f_n} \right\} \right|
&\leq \| p_{f_n} - p_0 \|_{\mu} \cdot \left\| \phi_{f_n} - \tilde{\phi}_{f_n} \right\|_{\mu} \notag \\
&\leq \| p_{f_n} - p_0 \|_{\mu} \cdot \left\|\phi_{f_n} - \tilde{\phi}_{f_n} \right\|_{\mu} \notag \\
&= O_P\left( n^{-\frac{k+1}{2k+3}} \right) \cdot O^+\left( J_n^{-(k+1)} \right) \notag \\
&= O_P^+\left( n^{-\frac{k+1}{2k+3}} n^{-\frac{k+1}{2k+3}} \right) \notag \\
&= o_P\left( n^{-1/2} \right), \label{eq:mean-difference-bound}
\end{align*}

Thus, we here finished proving Theorem~\ref{thm:Efficient Influence Curve Approximation}.
\end{proof}

\subsection{Proof of Theorem~\ref{thm:plugin_HAL_asymptotic_efficiency}}
\begin{proof}
Given we have established $P_n D^*_{f_n}=o_P(n^{-1/2})$, we have
\[
\Psi^F(f_n)-\Psi^F(f_0)=(P_n-P_0)D^*_{f_n}+R(P_{f_n},P_{f_0})+o_p(n^{-1/2}).\]
The Positivity Condition implies that the rate of $d_0(f_n,f_0)$ has the same rate of convergence in $L^2(\mu)$-norm. Then, by the Negligible Remainder Assumption in Theorem~\ref{thm:plugin_HAL_asymptotic_efficiency}, we can generally bound $R(P_{f_n},P_{f_0})$ by $\| p_{f_n}-p_{f_0}\|^2_{\mu}$.
So then $R(P_{f_n},P_{f_0})=O_P(n^{-\frac{2k+2}{2k+3}})$.
We also have that $\{D_f: f\in D^{(k)}_U([0,1])\}$ is a $P_0$-Donsker class and we assume 
$P_0\{D^*_{f_n}-D^*_{f_0}\}^2\rightarrow_p 0$. 

Therefore,
\[
(P_n - P_0)\{D^*_{f_0} - D^*_{f_n}\}
= n^{-1/2}\,o_p(1)
= o_p(n^{-1/2}).
\]

This then proves 
\[
\Psi^F(f_n)-\Psi^F(f_0)= (P_n - P_0)D^*_{f_0}+o_P(n^{-1/2}).\]
This proves asymptotic efficiency of $\Psi(f_n)$ as an estimator of $\Psi(f_0)$. 
\end{proof}

\subsection{Verification of the Assumptions for Median}
\begin{proof}
Recall that for the median functional $\psi(f) = F^{-1}(0.5)$, the efficient influence function is given by
\[
D_f(X) = \frac{1}{f(\psi)} \bigl( \tfrac{1}{2} - I(X \le \psi) \bigr),
\]
where $f$ is the corresponding density and $\psi = \psi(f)$ is the median. Let $\psi_n = \psi(f_n)$ and $\psi_0 = \psi(f_0)$.
\paragraph{Verification of the Second Assumption}
For the second assumption, we aim to show that 
\[
P_0\{D_{f_n} - D_{f_0}\}^2 = o_p(1).
\]

We have $\Vert f_n - f_0\Vert_{\infty} = O_P^+(n^{-\frac{k+1}{2k+3}})$. Then the CDFs $F_n$ and $F_0$ are close to each other in the sense that $\Vert F_n - F_0\Vert_{\infty} \leq \Vert f_n - f_0\Vert_{1} \leq \Vert f_n - f_0\Vert_{\infty} = O_P^+(n^{-\frac{k+1}{2k+3}})$.

Since $F_n(\psi_n) = F_0(\psi_0) = \tfrac{1}{2}$, we have
\[
0 = F_n(\psi_n) - F_0(\psi_0)
  = \{F_0(\psi_n) - F_0(\psi_0)\} + \{F_n(\psi_n) - F_0(\psi_n)\}.
\]
By the mean value theorem, there exists some $\xi_n$ between $\psi_n$ and $\psi_0$ such that
\[
F_0(\psi_n) - F_0(\psi_0) = f_0(\xi_n)\,(\psi_n - \psi_0).
\]
Taking absolute values yields
\[
|f_0(\xi_n)|\,|\psi_n - \psi_0| = |F_n(\psi_n) - F_0(\psi_n)|.
\]
Since $|F_n(\psi_n) - F_0(\psi_n)| \le \|F_n - F_0\|_\infty = O_P^+\!\left(n^{-\frac{k+1}{2k+3}}\right)$
and $\|f_0\|_\infty = O(1)$, we obtain
\[
|\psi_n - \psi_0|
= O_P^+\!\left(n^{-\frac{k+1}{2k+3}}\right).
\]

We start by expanding
\begin{align*}
D_{f_n} - D_{f_0}
&= \frac{1}{f_n(\psi_n)} \bigl( \tfrac{1}{2} - I(X \le \psi_n) \bigr)
 - \frac{1}{f_0(\psi_0)} \bigl( \tfrac{1}{2} - I(X \le \psi_0) \bigr) \\
&= \frac{f_0(\psi_0) - f_n(\psi_n)}{f_n(\psi_n) f_0(\psi_0)} \bigl( \tfrac{1}{2} - I(X \le \psi_0) \bigr)
  + \frac{1}{f_n(\psi_n)} \bigl(I(X \le \psi_0) - I(X \le \psi_n)\bigr).
\end{align*}
If we bound the $L^2(P_0)$-norm for the two terms on the right hand side, then we bound the left hand side.

For the first term, write
\[
f_0(\psi_0) - f_n(\psi_n) = \bigl(f_0(\psi_0) - f_n(\psi_0)\bigr)
   + \bigl(f_n(\psi_0) - f_n(\psi_n)\bigr).
\]
Assuming $f$ is differentiable with bounded derivative in a neighborhood of $\psi_0$, the mean-value theorem gives
\[
f_n(\psi_0) - f_n(\psi_n) = - f_n^{(1)}(\xi_n)(\psi_n - \psi_0)
\]
for some $\xi_n$ between $\psi_n$ and $\psi_0$. The $L^2(P_0)$-norm of the first term is bounded by $\Vert f_n - f_0\Vert_{P_0} + \vert\psi_n - \psi_0\vert = O_P^+(n^{-\frac{k+1}{2k+3}}) = o_P(1)$.

For the second term, note that WLOG, we have $\psi_0 < \psi_n$. Since $\psi_n-\psi_0$ converges to zero we have that the $L^2(P_0)$-norm of the latter indicator converges to zero.

\paragraph{Verification of the Third Assumption}
We aim to show that
\[
R(P_{f_n},P_0) = o_p(n^{-1/2}).
\]
\begin{align*}
R(P_{f_n},P_0)
&= \psi_n - \psi_0 + P_0 D_{f_n} \\
&= \psi_n - \psi_0 + \frac{F_0(\psi_0) - F_0(\psi_n)}{f_n(\psi_n)} \\
&= \psi_n - \psi_0 + \frac{f_0(\tilde{\xi}_n)}{f_n(\psi_n)} (\psi_0 - \psi_n) \\
&= \frac{f_n(\psi_n) - f_0(\tilde{\xi}_n)}{f_n(\psi_n)} (\psi_n - \psi_0) \\
& \leq \Vert f_n - f_0\Vert_{\infty} \cdot \vert\psi_n - \psi_0\vert = O_P^+(n^{-\frac{k+1}{2k+3}}) \cdot O_P^+(n^{-\frac{k+1}{2k+3}}) = o_P(n^{-1/2}).
\end{align*}

\end{proof}



\subsection{Proof of Theorem~\ref{thm:single-step-hal-tmle}}
For a single-step HAL-TMLE, knowing that $f_n^1$ converges to $f_0$ at a faster rate than $n^{-1/4}$, we have $P_n D^*_{f_n^1}= o_p(n^{-1/2})$, according to \citet{van2017generally}.
And thus the same proof of asymptotic efficiency holds for HAL-TMLE as well by extending all the $f_n$ to $f_n^1$. This is why we claim the proof above to be a standard TMLE proof.